\chapter*{Záver}
\addcontentsline{toc}{chapter}{Záver}
\markboth{Záver}{Záver} 

V súčastnosti, kedy je svet digitálne prepojený takmer na každej úrovni, spoločnosti spoločnosť pre svoje fungovanie potrebuje zaistiť primeraný (vysokú)
úroveň KIB. Chýbajúci odborníci KIB na Slovensku inšpirovali zadanie našej predošlej bakalárskej práce a neskôr aj tejto diplomovej práce. V zadaní sme
si dali za úlohu vytvoriť široko koncipovaný podporný systém na podporu analýzy rizík, cieľ presahujúci rozsah diplomovej práce a kapacity jedného človeka.

Konfrontovaní skutočným rozsahom a komplexitou, ktorú by si expertný systém vyžadoval, sme si položili reálny cieľ: navrhnúť a implementovať
malý systém s obmedzenou funkcionalitou kopírujúci funkcionalitu popísanú v kapitole \emph{\ref{chap:risk-analysis} -- \nameref{chap:risk-analysis}}.
Výsledný systém z kapitoly \emph{\ref{chap:application} -- \nameref{chap:application}} poslúžil ako východiskový bod, z ktorého sme v kapitole
\emph{\ref{chap:future-works} -- \nameref{chap:future-works}} načrtli funkcionalitu rozsiahlejšieho systému. Popísali sme problémy, ktoré ešte potrebuje
vyriešiť na to, aby sme mohli vyvinúť plnohodnotný expertný systém.

V diplomovej práci sme ukázali výsledok, ktorý sa nám počínajúc bakalárskou prácou podaril spraviť v priebehu takmer troch rokov. Čitateľovi
sme podali vedomosti, ktoré potrebuje, aby komplexne vyriešil KIB v organizácii. Analyzovali sme požiadavky informačného systému, ktorý by mu
pomohol zaviesť ISMS a vykonať analýzu rizík v organizácií. Získali sme predstavu, aká by bola zložitosť návrhu, implementácie, údržby a
prevádzkovania takého systému.

Predložený systém je \emph{proof-of-concept} a ešte by si vyžadoval praktický projekt, na ktorom by sme mohli objektívne posúdiť jeho prínos pre
prax a opodstatnenosť plnohodnotného softvérového riešenia.
